\documentclass[12pt, a4paper]{article}

\usepackage[utf8]{inputenc}
\usepackage[T1]{fontenc}
\usepackage[brazil]{babel}
\usepackage[top=3cm, bottom=2cm, left=3cm, right=2cm]{geometry}
\usepackage{setspace}
\usepackage{enumitem}
\usepackage{parskip}
\usepackage{ragged2e}

\begin{document}

\begin{flushright}
    Curitiba, 05 de março de 2026.
\end{flushright}

\vspace{0.5cm}

\noindent\textbf{Disciplina:} Experiência Criativa: Implementando Sistemas com Criptografia.

\noindent\textbf{Professores:} Altair Olivo Santin / Aramis Hornung Moraes

\vspace{0.3cm}

\noindent\textbf{Estudantes:}

\begin{itemize}[leftmargin=2cm, label={}]
    \item Gerard Gonzalez
    \item Gustavo Batista de Souza
    \item Henrique Bernardo Maciel
    \item Kauã Garcia Reschetti Rubbo
    \item Wellington Breno Evangelista Costa
\end{itemize}

\vspace{0.8cm}

\begin{center}
    \textbf{\large Portal para Registro de Propriedade}
\end{center}

\vspace{0.5cm}

\noindent Esse projeto tem o objetivo de fazer o cadastro de proprietário de qualquer produto que contenha número de série único, cujo objetivo é auxiliar o proprietário, a polícia e outros cidadãos a encontrar o proprietário.

\vspace{0.6cm}

\noindent\textbf{Páginas:}

\begin{itemize}
    \item Home
    \item Cadastro de usuário
    \item Cadastro de produto
    \item MFA
    \item Recuperação de senha
    \item Confirmação de cadastro
\end{itemize}

\end{document}
